\section{Introduction}


\subsection{Motivation}
The notion of \emph{ambient isotopy} is used in classical knot theory
to distinguish between different knot classes. It follows from a result
of Fisher \cite[Theorem~15]{fisher_1960} that 
two knots $\g_1,\g_2\in C^0(\R/\Z,\R^3)$
are ambiently isotopic if there exists an orientation preserving
homeomorphism $h:\R^3\to\R^3$ such that $h\circ\g_1=\g_2$. In concrete applications
it might be hard to come up with such a homeomorphism, but in the smooth category one
has the following useful stability theorem which to the best of our knowledge
has first been explicitly stated and proven by Diao et al.\@ who give also credit
to Ken Millet.
\begin{externaltheorem}[$C^1$-isotopy stability {\cite[Lemma 3$\cdot$2]{diao-etal_1999},
\cite{reiter_2005}, \cite[Proposition 3.1]{denne-sullivan_2008}}]
\label{thm:C1-stability}
For any immersed embedded curve $\g\in C^1(\R/\Z,\R^3)$ there is some number
$\epsilon_\g>0$ such that all curves $\eta\in C^1(\R/\Z,\R^3)$ with 
$\|\g-\eta\|_{C^1} < \epsilon_\g$ are embedded and ambiently isotopic to $\g$.
\end{externaltheorem}
This result also follows from the fact that any $C^1$-isotopy can be extended
to an ambient isotopy in arbitrary dimensions and co-dimensions which was proven\footnote{
For a proof in the simpler setting of the $C^2$-category see 
\cite[Chapter 8]{hirsch_1994}.}
by the first author in \cite{blatt_2009a}. The $C^1$-stability of isotopy is used
quite frequently in \emph{geometric knot theory} where one produces and investigates
specific knot representatives as minimizers or critical points of so-called
knot energies, 
to obtain more information about their knot classes; see, e.g., the
surveys 
\cite{strzelecki-vdm_2013a,strzelecki-etal_2013a,strzelecki-vdm_2014,blatt-reiter_2014,strzelecki-vdm_2018}  
and the many references therein. 
However, for the most interesting and mathematically more challenging
scale-invariant knot energies such as the M\"obius energy\footnote{In addition to 
scale-invariance the M\"obius energy is also invariant under M\"obius transformations
of the ambient Euclidean space \cite[Theorem~2.1]{freedman-etal_1994}, which lent
this energy that name.}
introduced by O'Hara \cite{ohara_1991a}
the underlying energy spaces do \emph{not} embed into $C^1$. Therefore,
one cannot expect, e.g., that minimizing sequences for these energies subconverge
in $C^1$, so that it is a priori not clear if the limit curves still belong to the
prescribed knot class. One of our goals in the present paper is therefore
to develop criteria
for stability of knot classes at low regularity.

For that we consider the slightly weaker notion
of \emph{knot equivalence} instead of ambient isotopy. Two knots $\g_1,\g_2\in
C^0(\R/\Z,\R^3)$ are said to be \emph{equivalent}, denoted by $\g_1\sim\g_2$, if and
only if there is a homeomorphism $h:\R^3\to\R^3$ such that $h(\g_1(\R/\Z))=\g_2(\R/\Z)$;
cf.\@ \cite[p. 4]{crowell-fox_1977}. It is known, e.g., that the trefoil knot is not
ambiently isotopic to its mirror image; see \cite[Theorem~3.39 (b)]{burde-zieschang_2014}. The corresponding reflection, however,
is an orientation reversing homeomorphism of $\R^3$ mapping the image of the trefoil
knot onto its mirror image, so that the two knots are equivalent.  So, there are
more (ambient) isotopy types than knot equivalence classes in knot space. On the other
hand, if two tame
arclength parametrized knots are equivalent, then one
can show that they are either ambiently
isotopic, or one is ambiently isotopic to the other one's mirror image. 

Gordon and Luecke studied in \cite{gordon-luecke_1989}
the fundamental groups of the knots' complements, also 
called \emph{knot groups}, and their peripheral subgroups. Combining their work
with results of Waldhausen \cite{waldhausen_1968} yields the following sufficient
criteria for knot equivalence.
\begin{externaltheorem}[Knot equivalence via knot complements
{\cite{waldhausen_1968,gordon-luecke_1989}}]\label{thm:gordon-luecke}
Let  $\g_1,\g_2\in C^0(\R/\Z,\R^3)$ be two tame knots.
\begin{enumerate}
\item[\rm (i)] If their complements $\R^3\setminus\g_i(\R/\Z)$, $i=1,2$, are homeomorphic,
then $\g_1\sim\g_2$.
\item[\rm (ii)] If $\g_1,\g_2$ are prime and their knot groups 
$\pi_1(\R^3\setminus\g_i(\R/\Z))$, $i=1,2,$ are isomorphic,  then $\g_1\sim\g_2$. 
\item[\rm (iii)] If the knot groups $\pi_1(\R^3\setminus\g_i(\R/\Z))$, $i=1,2,$ 
are isomorphic and if this isomorphism induces isomorphic peripheral subgroups $P_{\g_i}$, $i=1,2,$
then $\g_1\sim\g_2$ even if the $\g_i$ are composite knots.
\end{enumerate}
\end{externaltheorem}


Recall that for a set $M\subset\R^n$ with path-connected complement $\R^n\setminus M$
two closed loops $c_1,c_2\in C^0(\R/\Z,\R^n\setminus M)$ are equivalent if they are
homotopic, and that the fundamental group $\pi_1(\R^n\setminus M)$ consists of
the  respective equivalence classes. The group  structure is induced by concatenation
of closed loops containing a fixed base point in $\R^n\setminus M$. The peripheral
subgroup $P_M\subset\pi_1(\R^n\setminus M)$ then consists of all equivalence classes that
possess representative loops arbitrarily close to $M$.


\subsection{Isomorphic fundamental groups via localized distortion}
In order to apply the results on knot equivalence, i.e., Theorem~\ref{thm:gordon-luecke},
under low regularity assumptions
we are going to use a localized version of Gromov's distortion \cite{gromov_1978}. 
For a closed, path-connected set $M\subset\R^n$ its \emph{distortion} is defined as
the quantity
\begin{equation}\label{eq:distortion}
\delta(M):=\sup_{x,y\in M\atop x\not= y}\frac{d_M(x,y)}{|x-y|}
\qquad\ge 1,
\end{equation}
where $d_M(x,y)$ denotes the  
\emph{intrinsic distance} between the points $x,y$, i.e.,
the length of the shortest curve  in $M$ connecting $x$ with $y$. It is well-known that
the circle $c_{2\pi}$ satisfies $\delta(c_{2\pi})=\frac{\pi}2$ uniquely minimizing 
the distortion among all closed curves (see \cite[Proposition 2.1]{kusner-sullivan_1997} 
for a nice and short proof of that fact).  
Moreover, any non-trivial tame
knot $\g\in C^0(\R/\Z,\R^3)$ has distortion at least
$\frac{5\pi}3$, which was shown by Denne and Sullivan in \cite{denne-sullivan_2009}.
So, distortion detects knottedness, but it is not useful to distinguish between different
knots. For example, distortion cannot prevent the pull-tight phenomenon 
where knotted subarcs of a curve 
can be scaled down to vanish in the limit without any effect
on the curve's distortion \cite{mullikin_2007}. In addition, O'Hara constructed infinitely many
different prime knots with a uniform upper bound on their distortion 
\cite[Theorem~3.5]{ohara_1992a}.
On the other hand, Gromov asked if there is a universal bound on the distortion 
under which \emph{every}
knot type can be represented.
This was answered to the negative by Pardon in
\cite[Theorems 1.1 \& 1.2]{pardon_2011}
who also points out~\cite[p.~638]{pardon_2011} that there are 
even wild knots with finite distortion such as \cite[Figure~5]{crowell-fox_1977}.

Nevertheless, a universal bound on a
localized version of distortion
together with a comparable $L^\infty$-distance guarantees knot equivalence as
we will see in Corollary \ref{cor:equiv-knots} below.
\begin{definition}\label{def:local-distortion}
For a nonempty closed, path-connected set $M\subset\R^n$ and $r>0$ we define
the \emph{local distortion of $M$ at scale $r$} as
\begin{equation}\label{eq:local-distortion}
\delta(M,r):=\sup_{x,y\in M\atop 0<|x-y|\le 2r}\frac{d_M(x,y)}{|x-y|}.
\end{equation}
We say that \emph{$M$ has local distortion strictly below $\varepsilon>0$}
if there is some $r_M>0$ such that $\delta(M,r_M)<\varepsilon$,
and we call any such $r_M>0$ an \emph{admissible scale}.
%The set
%$M$ is said to have \emph{sub-quarter-circular distortion (SQC)} if \phso{there exists a scale $r_M>0$
%such that $\delta(M,r_M)<$} \philipp{has local distortion strictly below $\varepsilon = \frac{\pi}{\sqrt{8}}$}.
\end{definition}

Our first result on fundamental groups is stated for general 
subsets of $\R^n$ and might therefore
be of independent interest. Its formulation uses the
dimension-dependent constant
\begin{equation}\label{eq:dekster-jung}
g_n:=\sqrt{\frac{2n-2}n}\arcsin\sqrt{\frac{n}{2n-2}}
\quad\Fo n\ge 3,
\end{equation}
converging in a strictly decreasing manner
to $g_\infty:=\frac{\pi}{\sqrt{8}}$ as $n\to\infty$. Notice 
that
$g_3= \frac2{\sqrt{3}}\arcsin\frac{\sqrt{3}}2
= \frac{\sqrt{32}}{\sqrt{27}}g_\infty$ equals the
distortion of one third of a circle, whereas $g_\infty$
is the distortion of a quarter circle.
\begin{theorem}[Isomorphic fundamental groups]\label{thm:isomorphic}
Let $n\ge 3$.
If two nonempty closed, path-connected sets $M_i\subset\R^n$ 
with path-connected complements
$\R^n\setminus M_i$ have local distortion strictly
below $g_n$
at scales $r_{M_i}>0$ for $i=1,2$, and satisfy
\begin{equation}\label{eq:hausdorff}
\dist_\HH(M_1,M_2)<\frac14 \min\{r_{M_1},r_{M_2}\},
\end{equation}
then there exists an isomorphism $J:\pi_1(\R^n\setminus M_1)\to
\pi_1(\R^n\setminus M_2)$ such that its restriction to the peripheral
subgroup $P_{M_1}\subset \pi_1(\R^n\setminus M_1)$ maps $P_{M_1}$ isomorphically onto 
$P_{M_2}\subset \pi_1(\R^n\setminus M_2)$.
\end{theorem}


\subsection{Stability of knot equivalence}
As an immediate consequence of Theorem~\ref{thm:isomorphic} in combination with
part (iii) of Theorem~\ref{thm:gordon-luecke} we obtain a simple
sufficient geometric criterion for knot equivalence.
\begin{corollary}[Knot equivalence via local distortion]
\label{cor:equiv-knots}
Two tame knots $\g_i\in C^0(\R/\Z,\R^3)$ with local 
distortion 
 $\delta(\g_i(\R/\Z),r_{\g_i})<g_3$
at scales $r_{\g_i}>0$
for $i=1,2$, and
$\dist_{\HH}(\g_1(\R/\Z),\g_2(\R/\Z))$
$<\frac 14 \min\{r_{\g_1},r_{\g_2}\}$, are equivalent.
\end{corollary}

Notice that a small Hausdorff-distance, even an arbitrarily
small $L^\infty$-distance $\|\g_1-\g_2\|_{L^\infty}$
\emph{alone} does not guarantee knot equivalence:
One can insert small knotted arcs at any scale to produce highly complex
knots arbitrarily close to the round circle. Essential 
is the combination of controlled local
distortion at some scales and  
the Hausdorff-distance bounded in terms of these scales.
With Corollary \ref{cor:equiv-knots} we can prove 
the following stability of knot equivalence in the 
Lipschitz category.


\begin{corollary}[Lipschitz stability of knot equivalence]\label{cor:lipschitz-stability}
Let $\g:\R/\Z\to\R^3$ be an arclength parametrized tame 
knot with 
local distortion $\delta(\g(\R/\Z),r_\g)< g_3$
at scale $r_\g>0$. Then for all $h\in C^{0,1}(\R/\Z,\R^3)$
there exists $\epsilon_{\g,h}>0$ such that $\g_\epsilon:=\g+\epsilon h$ is
an immersed knot equivalent to $\g$ for all $|\epsilon| \le \epsilon_{\g,h}$.
\end{corollary}
Note that Lipschitz continuous curves are differentiable a.e., and we call such curves
\emph{immersed} if $v_\g:=\essinf_{u\in\R/\Z}|\g'(u)|>0$.

For $C^1$-knots $\g$ one has, of course, $\delta(\g(\R/\Z),r)\to 1 $ as $r\to 0$, so that
one finds a scale $r_\g$ sufficiently small so that $\g$ has 
local distortion strictly below $g_3$.
In fact, this observation implies a weaker version of Theorem~\ref{thm:C1-stability} where ambient isotopy is replaced by knot equivalence.
But we will prove here
that a fractional Sobolev regularity (below $C^1$) suffices, 
which leads to the following fractional stability result for knot equivalence.
\begin{corollary}[Fractional stability of knot equivalence]
\label{cor:fractional-stability}
For any arclength parametrized knot $\g\in W^{\frac32,2}(\R/\Z,\R^3)$ there exists
$\epsilon_\g>0$ such that every arclength parametrized
curve $\eta\in B_{\epsilon_\g}(\g)\subset
W^{\frac32,2}(\R/\Z,\R^3)$ is a knot equivalent to $\g$.
\end{corollary}
The preceding statement can be shown to hold true also
if $\g$ and
$\eta$ are only immersed and not necessarily arclength
parametrized, but we will not discuss the details here.
Open at this point, however, is if the bound on local 
distortion in Corollaries \ref{cor:equiv-knots} and 
\ref{cor:lipschitz-stability} already
implies tameness of the knots, or some improved regularity. 


\subsection{Symmetric critical points for scale-invariant knot energies}
The fractional Sobolev space $W^{\frac32, 2}(\R/\Z,\R^3)$ is the natural
energy space for  the 
\emph{M\"obius energy}
\begin{equation}\label{eq:emob}
E_\textnormal{M\"ob}(\g):=\int_{\R/\Z}\int_{\R/\Z}\left(
\frac1{|\g(s)-\g(t)|^2}-
\frac{1}{d_\g(\g(s),\g(t))^2}\right)|\g'(s)||\g'(t)|\,dsdt.
\end{equation}
Indeed, an arclength parametrized
curve $\g:\R/\Z\to\R^3$ has finite M\"obius energy  
%$\emob$ or $\TP_2$ 
if and only if 
$\g$ is embedded and of
class $W^{\frac32,2}(\R/\Z,\R^3)$ as shown in \cite{blatt_2012a} 
(see 
\cite[Theorem~3.2 (i)]{gilsbach-vdm_2018} for an explicit a priori bound
on the fractional seminorm in terms of  the energy). 
Freedman et al.\@ managed to
use the M\"obius invariance of $\emob$ to find minimizing knots in any given prime
knot class \cite[Theorem~4.3]{freedman-etal_1994}.
Critical points of the M\"obius energy were
shown to be smooth and even real-analytic 
\cite{blatt-etal_2016,blatt-vorderobermeier_2019} 
without using the invariance of $\emob$ under 
M\"obius transformations. 
But how can we find critical knots  possibly different from the absolute minimizers? 
One option is to use symmetry as carried out in \cite{gilsbach-vdm_2018} for 
O'Hara's
energies above scale-invariance whose energy spaces \emph{do} 
embed into $C^1$.

As in that 
paper we consider here for the M\"obius energy a discrete rotational symmetry. To be
more precise, for $p\in\N$, $p\ge 2$, we call a 
curve $\g\in C^0(\R/\Z,\R^3)$
\emph{($p$-rotationally) symmetric} if and only if
\begin{equation}
\label{eq:rot-symm}
\textstyle
\Rot_{\frac{2\pi}p}\circ\g(s -\frac1{p})=\g(s)\quad\Foa s\in\R/\Z, 
\end{equation}
where $\Rot_{\beta}\in SO(3)$ denotes the rotation about a fixed axis, say,
the $e_3$-coordinate axis, with rotation angle $\beta$.
Define for a given knot equivalence class $\mathcal{K}$ the 
\emph{($p$-rotationally) symmetric
subset}
\begin{equation}\label{eq:symmetric-subset}
\Sigma_p(\mathcal{K}):=\{\eta\in C^{0,1}(\R/\Z,\R^3):|\eta'|>0
\,\,\textnormal{a.e.},\, [\eta]=\mathcal{K},\,
\eta\,\,\textnormal{satisfies}\,\eqref{eq:rot-symm}\}.
\end{equation}
\begin{theorem}[Symmetric critical prime knots]
\label{thm:symm-crit-prime}
If the symmetric subset $\Sigma_p(\mathcal{K})$
of a prime knot equivalence class $\mathcal{K}$ contains a knot with
finite M\"obius energy or, equivalently, a 
$W^{3/2,2}$-knot,
then
there exists an arclength parametrized knot $\g\in C^\omega(\R/\Z,\R^3)
\cap\Sigma_p(\mathcal{K})$
such that
\begin{equation}\label{eq:minimizing}
\emob(\g)=\inf_{\Sigma_p(\mathcal{K})}\emob(\cdot),
\end{equation}
and  $\g$ is a critical knot, i.e., $d\emob (\g)[h]=0$ for all
$h\in W^{\frac32,2}(\R/\Z,\R^3)$. Moreover,
all torus knot classes $\mathcal{T}(a,b)$ for co-prime $a,b\in\Z\setminus\{0,\pm 1\}$
contain at least two different critical knots for the M\"obius energy.
\end{theorem}

One may prove a similar result for other scale-invariant knot energies
with energy spaces $W^{1+1/p,p}$
such as those from O'Hara's family; cf.\@ \cite{ohara_1992a,
blatt-etal_2022b}.

It might be easier to check the symmetries for polygonal
curves, so instead of a knot with finite M\"obius energy
we might require in Theorem \ref{thm:symm-crit-prime} that
$\Sigma_p(\mathcal{K})$ contains a polygonal knot.
Indeed, 
smoothing the corners such that the symmetry is preserved 
yields an immersed embedded smooth
curve in $\Sigma_p(\mathcal{K})$
which has finite Möbius energy due 
to~\cite{blatt_2012a} and the energy's invariance under reparametrization. We do not know, on the other hand, if it suffices
to assume in Theorem \ref{thm:symm-crit-prime} only
that $\mathcal{K}$ is tame and that
$\Sigma_p(\mathcal{K})$ is nonempty.



We mentioned before that two equivalent arclength parametrized tame
knots $\g_1,\g_2:\R/\Z\to\R^3$ are either already ambiently isotopic, or
$\g_1$ is ambiently isotopic to the mirror image $\g_2^*$ of $\g_2$.
The $O(3)$-invariance of the M\"obius energy therefore
allows us to 
strengthen Theorem~\ref{thm:symm-crit-prime} to ambient isotopy
classes.
\begin{corollary}[Symmetric critical prime knots in ambient isotopy
classes] \label{cor:symm-crit-prime}
Let $\mathcal{K}$ be an ambient isotopy class of prime knots and $\Sigma_p(\mathcal{K})$ be the class of $p$-rotationally symmetric regular Lipschitz
curves whose ambient isotopy class equals $\mathcal{K}$. If
$\Sigma_p(\mathcal{K})$ contains a knot of
finite M\"obius energy, then there is an arclength parametrized
minimizing real analytic knot in $\Sigma_p(\mathcal{K})$ such that
\eqref{eq:minimizing} holds, and this knot is $\emob$-critical. If,
more specifically,
$\mathcal{K}$ equals the ambient isotopy class of torus knots,
$\mathcal{T}(a,b)$ for co-prime $a,b\in\Z\setminus\{0,\pm 1\}$, then there
are at least two such $\emob$-critical knots.
\end{corollary}
Kim and Kusner use in \cite{kim-kusner_1993}
symmetric criticality to construct $\emob$-critical torus 
knots (even with explicit formulas for their energy values). 
Their numerical experiments suggest, that
for (co-prime) parameters $a,b$ with $|a|,|b|>2$ 
these critical
points seize to be stable. Analytically, this remains an
open question. 




\subsection{Strategy and outline of the paper}\label{sec:strategy}
In Section~\ref{sec:distortion} we analyze the effects
of the smallness
condition on local distortion  on general
closed, path-connected subsets $M\subset\R^n$. For 
each $x$ near (but not in) $M$ we find
a direction that ``points away from $M$''. These
directions are glued together in Section~\ref{sec:pseudo}
with a partition of unity to construct a kind of
pseudo-gradient vector field for the distance function $\dist(\cdot,M)$.  We solve
uniquely and globally two different flow equations with respect
to this vector field to define homotopies in the complement
$\R^n\setminus M$ of $M$. These mappings are then used in Section~\ref{sec:iso} 
for $M:=M_1$ and $M_2$ in combination with the sets' 
controlled local distortion
to construct  an
isomorphism between the fundamental groups
$\pi_1(\R^n\setminus M_1)$ and $\pi_1(\R^n\setminus M_2)$ as stated in 
Theorem~\ref{thm:isomorphic}. 

We use in Section~\ref{sec:stability} the immediate Corollary
\ref{cor:equiv-knots} to prove the Lipschitz stability and fractional stability
of knot equivalence, i.e., Corollaries \ref{cor:lipschitz-stability} and
\ref{cor:fractional-stability}. The first one relies on the fact that for curves the smallness condition on local 
distortion  
is stable with respect to Lipschitz perturbations. 
For
the fractional stability we recall in Corollary \ref{cor:seminorm-bilip}
that a small local fractional semi-norm of the
tangent implies a local bilipschitz estimate, which can be turned into a local
distortion bound for embeddings.

More analytical work is required for the existence
of symmetric $\emob$-critical prime knots as stated in
Theorem~\ref{thm:symm-crit-prime}. Since minimizing sequences converge only
weakly in $W^{\frac32,2}$ to a 
limit curve, we can \emph{not} relie on the (strong) 
$W^{\frac32,2}$-stability of Corollary \ref{cor:fractional-stability} to
guarantee the prescribed knot equivalence class in the limit. Nevertheless, 
the limiting
curve $\g$ has finite M\"obius energy because of the energy's lower-semicontinuity, 
so that we have
an arclength parametrized tame
$W^{\frac32,2}$-knot $\g$ as a limit. Exactly as in the proof of
Corollary \ref{cor:fractional-stability} we therefore find a scale  at which
$\g$ has its local distortion strictly
below the universal constant
$g_3$. The same holds true for all members of
the minimizing sequence, because each of them is an arclength parametrized
$W^{\frac32,2}$-knot as well. But the $L^\infty$-convergence of that sequence
towards $\g$ might be too slow to satisfy the assumptions of Corollary \ref{cor:equiv-knots}
for any sequence member.
In other words, the scales at which 
the members of the minimizing sequence have local 
distortion below $g_3$
might decay to zero much faster than their $L^\infty$-distance to $\g$. 
The analytic reason for that behaviour lies in the fact that we do \emph{not}
have uniform absolute continuity of the integrals defining the fractional
semi-norms of the tangents of the minimizing sequence. So, we are led to 
study possible concentrations in these semi-norms to cut them out, and replace
the corresponding subarcs of the curves by straight segments. Since such a substitution
generally destroys the $W^{\frac32,2}$-regularity\footnote{\label{foot:harmonic-substitution}For arclength parametrized embedded curves, $W^{\frac32,2}$-regularity is equivalent to finite Möbius energy according to~\cite{blatt_2012a}.
Due to~\cite[Corollary~1.3]{freedman-etal_1994}, such a curve has a local bilipschitz constant (cf.\@ Footnote~\ref{foot:bilip}), arbitrarily close to~$1$.
However, replacing a subarc of a curve 
by a line is likely to produce sharp corners at its endpoints
which results in a bilipschitz constant strictly larger 
than~$1$.},
a careful analytic
investigation of that
procedure based on some subtle BMO-type estimates
is required in order to keep the local
distortion of these modified curves under control. This analysis is carried out
in Section~\ref{sec:harmonic}.

In Section~\ref{sec:compactness} we interpret
the fractional semi-norms of sequences of $W^{\frac32,2}$-knots as Radon measures
on $\R/\Z\times\R/\Z$ to pass to a weak limiting  Radon measure in order to analyze
its set of concentration points.  At  a carefully chosen scale around these 
concentration points we apply the cutting procedure of 
Section~\ref{sec:harmonic}
to members of the minimizing sequence with sufficiently large
index. The result
is a weak fractional compactness theorem in the set of symmetric, arclength 
parametrized $W^{\frac32,2}$-knots (see Theorem~
\ref{thm:fractional-compactness}), which
might be of independent interest. This is the main tool in the existence proof
for symmetric $\emob$-minimizing knots in Section~\ref{sec:symm-critical}. A simple
symmetrization argument which  uses the invariance of the M\"obius energy under rotations
then leads to the fact that these minimizing knots are also critical points for
$\emob$, and therefore are real analytic. This
finalizes the proof of Theorem~\ref{thm:symm-crit-prime}.
